    \subsection{Framed self-linking of the swirl string and the spinorial selection of the lepton ladder}
        \label{subsec:framed_selflinking_spinorial}

        This subsection records the framed-tube self-linking of a closed torus-knot swirl
        string and the boundary class that selects its physical winding. It is the first
        genuinely non-circular topological result in the canon, in deliberate contrast to the
        calibrated $\alpha$-identities of the Formal Foundations (which re-express $\alpha$
        rather than predict it).

        \textbf{[DERIVED] Torus-surface self-linking.}
        Let $K=T(p,q)$ be a closed swirl string on the standard torus and let $U_{\rm tor}$ be
        the torus surface-normal framing. The framed self-linking is
        \begin{align}
            SL_{\rm tor}(T(p,q)) \;=\; pq,
            \label{eq:SL_torus_pq}
        \end{align}
        verified numerically to machine precision for $p\in\{2,3\}$ and $q=3,\dots,11$ with
        \emph{no} target injected. For the lepton family $T(2,q)$, $SL_{\rm tor}=2q$.

        \textbf{[DEFINITIONAL] Physical-framing winding.}
        For any single-valued physical framing $U_{\rm phys}$ of $K$ (the direction of the core
        swirl velocity $\vswirl$), the relative winding
        \begin{align}
            \chi \;:=\; SL(K,U_{\rm phys})-SL(K,U_{\rm tor})
                  \;=\; \frac{1}{2\pi}\oint_K \mathrm{d}\!\left[\angle_{U_{\rm tor}}(U_{\rm phys})\right]
                  \;\in\; \mathbb{Z}
            \label{eq:chi_def}
        \end{align}
        is an integer relative winding number, so $SL_{\rm phys}(T(p,q))=pq+\chi$.

        \textbf{[CRITICAL NOTE]} Any test that fixes the twist by $Tw=\mathrm{target}-Wr$ and
        then confirms the target is circular: it passes for \emph{any} integer target and
        proves only the Calugareanu identity $SL=Wr+Tw$. Only the target-free
        framing~\eqref{eq:SL_torus_pq} and the selection argument below are admissible.

        \textbf{[DERIVED NEGATIVE] Ordinary $U(1)$ closure gives $\chi=0$.}
        With positive longitudinal phase stiffness $A_\parallel>0$, the winding energy
        \begin{align}
            E(\chi)\;=\;E_0+\frac{2\pi^2 A_\parallel}{L}\,\chi^2
            \label{eq:E_chi}
        \end{align}
        is minimized over ordinary (periodic) $U(1)$ windings $\chi\in\mathbb{Z}$ at
        $\chi_{\rm ground}=0$, hence $SL_{\rm phys}=pq$ (even). Pure framed-curve geometry with
        positive stiffness therefore yields a \emph{boson} and cannot produce the lepton
        $pq+1$ ladder.

        \textbf{[CONDITIONAL DERIVED] Spinorial anti-periodicity gives $\chi=\pm1$.}
        If the core phase is a section of a spinor (rather than a $U(1)$) bundle, then under one
        circuit of $K$ the spin frame rotates by $2\pi$ and the phase is anti-periodic, so the
        admissible windings are $\chi\in 2\mathbb{Z}+1$. Minimizing~\eqref{eq:E_chi} over this
        sector gives the degenerate ground winding
        \begin{align}
            \chi_{\rm ground}=\pm1 \quad\Longrightarrow\quad
            SL_{\rm phys}(T(p,q))=pq\pm1,
            \qquad SL_{\rm phys}(T(2,q))=2q\pm1,
            \label{eq:SL_phys_pq1}
        \end{align}
        the two signs being the matter and antimatter sectors. The selection is a
        \emph{superselection} effect: the boson winding $\chi=0$ is always energetically lower
        (by $2\pi^2 A_\parallel/L>0$), so it is the boundary class, not the energetics, that
        renders the carrier fermionic. The result is independent of $(A_\parallel,L)$ for all
        positive values.

        \textbf{[OPEN] Derivation of the spinorial boundary condition.}
        The spinorial anti-periodicity of the core phase is not yet derived from the swirl-field
        dynamics; it is the single remaining premise of the $pq+1$ ladder.

        \textbf{[CONDITIONAL] Unification with the magnetic moment.}
        The spinor premise above is the \emph{same} posit as the spin-$\tfrac12$ assignment of
        the magnetic-moment sector (Route~B): one spinorial lift of the core phase yields both
        $g=2$ and $\chi=\pm1\Rightarrow pq\pm1$. If established, fermion statistics, the
        gyromagnetic ratio, and the lepton self-linking ladder collapse into a single mechanism.

        \textbf{[DERIVED] Falsifier.}
        An independently computed $T(2,q)$ Hopfion self-linking with $Q_H\neq 2q\pm1$ (e.g.\
        $T(2,5)\neq 11$) falsifies the conditional chain irrespective of the variational
        argument.
